\section{Data and Valuation Setup}

\subsection{Controlled layer separation}
The integrated experiment is organized as a controlled handoff across three modules. Team 8 governs only the gold-price process; Team 4 supplies the operating forecasts; and Team 5 maps those inputs through a common tax, WACC and terminal-value contract without a reconciled equity bridge. In the four-model comparison only the gold engine changes: the same Team 4 production and unit-cost forecasts, the same Team 5 corporate rules and the same WACC shocks are reused across runs.

\subsection{Valuation clock}
The valuation time is the 2 September 2026 close (20:00~UTC). Only flows after it are valued. Barrick's Q1--Q2 2026 actuals have already been realised and are excluded; they are used only to check levels. The quarter that contains the valuation date contributes its remaining \StubDays{} of 92 days, under the stated assumption that production and unit cost are uniform within the quarter. Eighteen quarters follow, to 31 December 2030 (\HorizonYears{} years). Revenue in each quarter uses the average simulated daily gold price over its remaining interval, and each flow is discounted from the midpoint of that interval.

\subsection{Calibration snapshot and carry}
The calibration uses the 2 September 2026 GLD call surface acquired through the IB API. After the common DTE $\geq75$ and numerical filters, the dense surface contains 605 eligible observations across 12 expiries and 146 strikes, with 79--653 days to expiry. The $8\times8$ Chebyshev--Chebyshev sample contains 64 distinct actual contracts spanning eight expiries and 20 strikes. The rate input is a same-date NSS fit to 14 U.S. Treasury tenors \citep{nss1994,treasury2026}, with a 2.0608~bp RMSE against the continuously compounded par-yield proxy; it is not a bootstrapped zero curve.

GLD's sponsor fee of 0.40\% a year reduces the gold backing of each share, so the GLD forward is $S e^{(r-q)T}$ with $q=0.40\%$ even though no cash is distributed \citep{spdr2025}. Two independent observations support this value: the GLD/gold log-ratio trend of $-0.396\%$ a year over 2022--2026, and COMEX gold futures, whose carry against the synchronous spot implies a GLD yield of \FuturesYieldLow--\FuturesYieldHigh\% at 0.5--0.8 years. The carry is applied exactly through $C(S,K;q)=e^{-qT}C(S,Ke^{qT};0)$.

The gold scenarios start from the XAUUSD spot of USD~4,391.32/oz at the valuation time, synchronous with the GLD close. The earlier draft started from Barrick's Q2 average realised price of USD~4,417/oz, which predates the valuation by two months.

\subsection{Probability measure}
Option parameters are estimated under $\mathbb Q$. The scenario drift is the integral-preserving forward curve of the Treasury fit. Read as a physical expectation, this assumes no gold risk premium, which is a common assumption for an asset with near-zero market beta but is not estimated here. The option-implied volatility shape is carried over as a conditional assumption; it is not a validated physical law. Expected flows are then discounted at WACC, as in an ordinary DCF. For comparison, Section~\ref{sec:results} also reports the explicit flows discounted at the risk-free curve. A full $\mathbb Q\rightarrow\mathbb P$ specification would need price, variance and jump risk premia (Section~\ref{sec:paper-physical-measure}).

\begin{table}[H]
\centering
\caption{Contract held fixed across gold engines.}
\label{tab:frozen_contract}
\begin{tabularx}{\columnwidth}{@{}>{\bfseries}p{0.25\columnwidth}X@{}}
\toprule
Block & Current experiment \\
\midrule
Gold layer & Black--Scholes/GBM, Heston, Bates--Poisson and Full Bates--Hawkes calibrated by Team 8 under $\mathbb Q$ with $q=0.40\%$. \\
Operations & Team 4 production and unit-cost forecasts from 2026Q3 (stub) to 2030Q4; actuals excluded. \\
Corporate mapping & Team 5 tax, growth--ROIC--reinvestment schedule, stochastic WACC and explicit terminal convention; no equity bridge. \\
Simulation & 8,192 paths, daily steps, one common WACC shock array. \\
\bottomrule
\end{tabularx}
\end{table}

\subsection{Option exercise convention}
GLD contracts are American-style \citep{cboe2026}. With $q=0.40\%$ below the Treasury rate, early exercise of a call is not optimal in the frictionless model. An 800-step CRR check on all 64 contracts gives a maximum American--European premium of USD~\AmericanPremiumMax; this tests the convention, not market frictions.

\subsection{Operating-model status}
Team 4's original operating study is richer than the object passed downstream. Cost of Sales is forecast from a log-space master chain; production is constructed bottom-up from ore processed, average grade and recovery. The valuation uses the resulting point forecasts, so mine-level operating uncertainty is not propagated jointly with gold-price uncertainty. Team 4's legacy price simulation and illustrative EBITDA/valuation output are excluded.

\subsection{Gold perimeter and reporting basis}
The Team 4 subset comprises Bulyanhulu, Carlin, Cortez, Kibali, Loulo--Gounkoto, North Mara, Pueblo Viejo and Turquoise Ridge. Its quarterly estimation window is 2019Q1--2025Q4, with North Mara beginning in 2019Q4. The operating quantities are already attributable: Nevada Gold Mines is 61.5\%, Pueblo Viejo 60\%, Kibali 45\%, Loulo--Gounkoto 80\%, and the Tanzanian mines 84\%. Kibali is equity-accounted, whereas the other interests are consolidated with partners' non-controlling interests \citep{barrickar2025}. Those partners must not be deducted again from attributable flows. Tanzania's contractual economic-benefit sharing still requires a separate cash-flow reconciliation.

The forecast covers the eight-mine subset, which produced 622 and 711~koz in Q1--Q2 2026 against company-wide totals of 719 and 796~koz. The proxy therefore omits the remaining gold mines. Cost of Sales includes depreciation and price-sensitive royalties; multiplying its sold-ounce cost by production gives a component-margin proxy, not EBITDA or reconciled EBIT. Loulo--Gounkoto's 2025Q4 accounting uplift also makes its cost anchor a normalization issue. Ownership, the North America IPO, cash and debt, share count and reserves must be fixed at the valuation date before any per-share result can be reported.
